\chapter{Теоретические основы и анализ проблематики}
\label{sec:teorie}

Эта глава посвящена теоретическим основам, на которых построено спроектированное устройство для тестирования и мониторинга. Понимание этих понятий необходимо для правильной интерпретации данных, полученных при анализе сети.

\section{Протоколы сетевого уровня и их различия в производительности}
\label{sec:ipv4_ipv6}

Современные компьютерные сети опираются на сосуществование протоколов IPv4 и IPv6 \cite{nastase2024}. Обе версии выполняют на сетевом уровне одну и ту же задачу — адресацию станций и доставку пакета между сетями, — но различаются настолько, что измерительный прибор должен обслуживать каждую из них отдельно. Одновременная работа обоих протоколов на одном интерфейсе, называемая двойным стеком, сегодня обычна в корпоративных сетях, и переносной анализатор должен уметь проверить работоспособность каждой ветви по отдельности.

Самое заметное различие — длина адреса, выросшая с 32 до 128~бит. Изменение адресного пространства, однако, не единственное вмешательство в конструкцию протокола. Заголовок IPv4 имеет переменную длину от 20 до 60~байт, поскольку за обязательной частью могут следовать дополнительные поля, и содержит, помимо прочего, собственную контрольную сумму и три поля для фрагментации. Заголовок IPv6, напротив, фиксирован и занимает 40~байт \cite{rfc8200}. Вся дополнительная функциональность перенесена в заголовки расширения, которые выстраиваются цепочкой за основным заголовком и на существование которых указывает поле Next Header.

У такой конструкции два следствия, которые проявляются именно при измерениях. Во-первых, маршрутизатор на пути пакета выполняет меньше работы: он не пересчитывает контрольную сумму, потому что её вообще убрали из заголовка, полагаясь на защиту канального и транспортного уровней, и не фрагментирует. Фрагментировать в IPv6 может только отправитель, который для этого должен заранее выяснить наименьший размер передаваемого блока на всём пути. С этим связано требование, чтобы каждая линия передавала пакет размером не менее 1280~байт. Во-вторых, IPv6 полностью отказался от широковещательной рассылки и заменил её групповой, что прямо отражается на том, как устройства в сегменте находят друг друга; этому посвящён подраздел \ref{sec:ipv6_adresace}.

С точки зрения измерения пропускной способности существенно, что вдвое больший основной заголовок — это накладные расходы, которые вычитаются из полезной нагрузки. При одинаковом размере кадра канального уровня на данные остаётся на двадцать байт меньше, поэтому измеренная прикладная пропускная способность у IPv6 получается чуть ниже. Как показывают сравнительные измерения Tuifaiga и соавторов \cite{tuifaiga2021}, IPv4 в локальных сетях действительно достигает несколько большей пропускной способности и меньшей задержки, чем IPv6, одинаково в Windows и Linux.

Различие, впрочем, нужно толковать осторожно. Его относительная величина зависит от того, какими блоками передаются данные: у длинных потоков с полноразмерными кадрами это единицы процентов, а у трафика из коротких пакетов может быть заметно больше. Поэтому измеренная разница между протоколами сама по себе говорит не о качестве сети, а о свойствах протокола, и при оценке нагрузочных тестов из главы \ref{sec:realizace} её нужно отличать от реальных неисправностей.

\section{Адресация и автоматическая настройка протокола IPv6}
\label{sec:ipv6_adresace}

Различие между протоколами не ограничивается размером заголовка. Для переносного анализатора важнее другой способ, которым станция вообще получает свой адрес, поскольку от него зависит, сможет ли прибор обратиться к другой стороне без всякой предварительной договорённости с её владельцем. Поэтому этому вопросу посвящён отдельный подраздел.

Адрес IPv6 имеет длину 128~бит, его структура определена архитектурой адресации \cite{rfc4291}, а сам протокол — документом \cite{rfc8200}. Запись делится на восемь групп по четыре шестнадцатеричные цифры, причём самый длинный непрерывный участок нулевых групп можно заменить двумя двоеточиями. Для тестового устройства важны три категории адресов. Адрес локального канала из диапазона \texttt{fe80::/10} возникает на каждом интерфейсе автоматически в момент поднятия линии и действует только в пределах одного физического сегмента. Глобальный адрес маршрутизируется во всём Интернете. Между ними стоят уникальные локальные адреса, которые для нашего устройства интереснее всего.

Уникальные локальные адреса задаются префиксом \texttt{fc00::/7}, причём восьмой бит различает способ выделения; при его установке, то есть для префикса \texttt{fd00::/8}, организация сама генерирует сорокабитный идентификатор псевдослучайным образом \cite{rfc4193}. Такой префикс с очень высокой вероятностью уникален, не маршрутизируется в Интернете и не зависит ни от какого провайдера. Именно эти свойства соответствуют природе переносного анализатора, который работает в закрытом тестовом сегменте, не имеет постоянного подключения к Интернету и чей адресный план не должен конфликтовать с адресами сети, куда прибор временно подключён.

Вторую половину адреса образует идентификатор интерфейса длиной 64~бита. Один из способов его получения, описанный в архитектуре адресации \cite{rfc4291}, исходит из физического адреса интерфейса и называется EUI-64. Сорокавосьмибитный адрес делится на две половины, между которыми вставляются два байта со значением \texttt{fffe}, а в первом байте инвертируется седьмой бит, различающий универсальное и локальное управление. Так физический адрес \texttt{00:e0:4c:68:02:16} даёт идентификатор \texttt{02e0:4cff:fe68:0216}. Для нашего устройства это практически важно: адрес другой стороны получается детерминированно, и его можно вычислить, а не прописывать жёстко в управляющем приложении.

У идентификатора, полученного из физического адреса, есть, однако, недостаток: он выдаёт оборудование станции и позволяет следить за ней, в какую бы сеть она ни подключилась. Поэтому современные операционные системы чаще используют идентификатор, который для данной сети постоянен, но из физического адреса не выводится \cite{rfc7217}, а рядом с ним ещё и временные адреса, которые регулярно меняются \cite{rfc8981}. Одна станция поэтому обычно несёт на одном интерфейсе несколько адресов сразу и, если так настроена, для исходящей связи предпочитает временный \cite{rfc8981}. Для измерительного прибора отсюда следует, что другая сторона может отвечать не с того адреса, к которому к ней обратились, — в практической части это оказалось существенным.

Обнаружением соседей на общей линии в IPv6 занимается механизм Neighbor Discovery \cite{rfc4861}, который заменяет протокол ARP из IPv4 и построен поверх ICMPv6. С его помощью устройства узнают физические адреса друг друга, проверяют достижимость соседей и находят маршрутизаторы в сегменте. Для нашего устройства ключевые два его сообщения. Сообщением Router Solicitation, отправляемым на групповой адрес всех маршрутизаторов, станция после поднятия интерфейса спрашивает, есть ли в сегменте маршрутизатор. Ответом служит сообщение Router Advertisement, которым маршрутизатор объявляет себя и вместе с тем префиксы, действующие в сегменте. Частью механизма является и обнаружение дублирующихся адресов, которым станция проверяет, что выбранный ею адрес в сегменте ещё не занят.

На сообщении Router Advertisement построена автоконфигурация адресов без сохранения состояния \cite{rfc4862}. Станция сама составляет полный адрес из объявленного префикса и собственного идентификатора интерфейса, без какого-либо сервера, который вёл бы учёт выдачи. Воспользуется ли она этой возможностью и как, определяют три флага в сообщении, значение которых сведено в таблице \ref{tab:priznaky}.

\begin{table}[htbp]
    \caption{Значение флагов в сообщении Router Advertisement}
    \label{tab:priznaky}
    \centering
    \begin{tabular}{clp{7.5cm}}
        \hline
        Флаг & Название & Значение для станции \\
        \hline
        M & Managed address configuration & адрес станция запрашивает у сервера DHCPv6 \\
        O & Other configuration & прочие данные, например адреса DNS, получает от сервера DHCPv6 \\
        A & Autonomous address configuration & из этого конкретного префикса станция может составить адрес сама \\
        \hline
    \end{tabular}
\end{table}

Принципиально, что флаги M и O относятся ко всему сообщению, а флаг A передаётся у каждого объявленного префикса отдельно. Маршрутизатор может в одном сообщении объявить несколько префиксов и для каждого решить, может ли станция вывести из него адрес самостоятельно. Оба пути не исключают друг друга и на практике обычно комбинируются, так что у станции может быть одновременно адрес, составленный без сохранения состояния, и адрес, выданный сервером.

Альтернативой с сохранением состояния является DHCPv6 \cite{rfc8415}, перенимающий роль, знакомую по IPv4: сервер ведёт учёт выданных адресов, привязывает их к идентификатору клиента и управляет сроком аренды. Ценой этого учёта он даёт администратору полный контроль над тем, какая станция какой адрес получит, и позволяет передать и другие параметры настройки.

Разница подходов для переносного тестового устройства существенна. Станция с IPv4 без работающего клиента DHCP адрес не получит никогда, потому что другого механизма протокол не предлагает. Ядро же операционной системы с IPv6 выдаёт адрес локального канала само, а при поднятом флаге A составит из объявленного префикса и адрес более широкой области действия, причём на станции не должна работать никакая служба настройки сети. Анализатор, правильно объявляющий префикс, может поэтому установить связь и с машиной, на которой ничего не настроено и над которой у оператора прибора нет никакого контроля. Проверка этого предположения на собранном прототипе описана в подразделе \ref{sec:overeni_autokonfigurace}.

\section{Генерирование сетевого трафика}
\label{sec:generovani}

Работу с сетевым трафиком можно разделить на два принципиально разных подхода. Пассивный мониторинг только наблюдает трафик и ничего в сеть не отправляет, а активное измерение само создаёт трафик и смотрит, как с ним справится сеть. Переносной анализатор должен уметь и то, и другое, потому что они отвечают на разные вопросы: пассивное наблюдение показывает, что в сегменте происходит на самом деле, активное измерение — на что сегмент способен. Этот раздел посвящён активной стороне, то есть генерированию трафика, следующий — пассивной.

Смысл генерирования в том, что в сеть отправляют трафик с известными свойствами и на приёмной стороне выясняют, что из него и как дошло. Из сравнения отправленного и принятого определяются основные параметры сети. Пропускная способность показывает, сколько данных в секунду сеть передаст. Потери показывают, какая доля отправленных пакетов не дошла. Задержку чаще всего измеряют как время оборота, то есть время между отправкой сообщения Echo Request протокола ICMP и приходом ответа, как это делает утилита ping. Наконец, джиттер (разброс задержки) описывает, насколько колеблется время передачи последовательных пакетов.

Генераторы трафика по уровню, на котором они работают, делятся на две группы. Генераторы потоков создают непрерывную передачу данных между клиентом и сервером и оценивают результаты на обеих сторонах соединения. Типичный представитель — инструмент iperf3 \cite{iperf3}. Выбор транспортного протокола определяет, что именно измеряется. Измерение поверх TCP даёт пропускную способность, достижимую для приложения, с учётом управления потоком и восстановления после потерь, поэтому результат отражает поведение реальной передачи данных. Измерение поверх UDP, наоборот, отправляет с заранее заданной скоростью без оглядки на состояние сети и кроме пропускной способности говорит о потерях и джиттере. Различие между передачей с установлением соединения и без него на высоких скоростях подробно разбирают Gál и соавторы \cite{gal2021}; их работа служит методической основой для измерений в этой работе. Для содержательного результата тест должен длиться достаточно долго, чтобы управление потоком успело установиться, и измерение нужно повторять.

Вторую группу составляют генераторы пакетов, которые собирают отдельные кадры выбранного типа слой за слоем. Они служат не столько для измерения пропускной способности, сколько для проверки, пройдёт ли определённый вид трафика через сегмент и как на него отреагируют остальные устройства. В языке Python для этого есть библиотека Scapy, которая позволяет сложить кадр из отдельных заголовков, отправить его и обработать ответы \cite{carter2024}. В коммерческих тестерах для той же цели стоит специализированная аппаратура, способная отправлять точно рассчитанные по времени кадры на полной скорости линии; программный генератор на обычном компьютере ограничен производительностью процессора и точностью таймеров операционной системы.

Скорость генерируемого трафика можно выразить двумя способами — в битах в секунду и в кадрах в секунду, — и эти величины не взаимозаменяемы. Обработка каждого кадра требует и в сетевом устройстве, и в компьютере определённой постоянной работы независимо от его размера. Поэтому на коротких кадрах устройство обычно исчерпывает свою производительность раньше, чем заполнит пропускную способность линии, и его предел задаётся числом кадров в секунду; на длинных кадрах, наоборот, раньше упирается в пропускную способность линии или шины. Методика измерения производительности сетевых устройств \cite{rfc2544} поэтому предписывает повторять измерение для ряда размеров кадра, для Ethernet это 64, 128, 256, 512, 1024, 1280 и 1518~байт, и пропускной способностью считает наибольшую скорость, при которой не теряется ни один кадр. Для IPv6 рекомендуются те же размеры кадра \cite{rfc5180}, но с оговоркой, что заголовок IPv6 на двадцать байт длиннее и самые короткие кадры могут не вместить все заголовки и данные измерительного инструмента.

Джиттер для UDP обычно оценивают способом, заимствованным из протокола RTP \cite{rfc3550}. Для каждой пары последовательно принятых пакетов определяют разность $D$ между разностью времён их прихода и разностью времён их отправки. Из этих разностей непрерывно вычисляется сглаженная оценка
\begin{equation}
    J_i = J_{i-1} + \frac{|D_{i-1,i}| - J_{i-1}}{16}.
    \label{eq:jitter}
\end{equation}
Поскольку используются только разности времён, часы отправителя и получателя не обязаны быть синхронизированы. По той же формуле считает джиттер и iperf3 \cite{iperf3}. Для толкования результатов полезно также различать среднюю и мгновенную скорость: программный генератор может выдерживать заданную среднюю скорость и при этом отправлять короткими пачками на полной скорости линии, которые принимающее устройство с небольшим буфером может не успеть обработать.

Активное измерение нагружает сеть, поэтому важно, где оно проводится. Методы \cite{rfc2544} разрабатывались для измерения отдельных устройств в изолированной лабораторной среде. Дополняющий документ \cite{rfc6815} прямо предостерегает от их применения в рабочей сети: тестовый трафик может перегрузить ресурсы, общие с пользователями, а результаты к тому же будут неверными, потому что потери от чужого трафика не отличить от потерь измеряемого устройства. Переносной прибор поэтому должен ограничивать нагрузочные тесты закрытым тестовым сегментом.

\section{Мониторинг и аудит безопасности сетевого трафика}
\label{sec:monitorovani}

Пассивный захват в системах на базе Linux построен на библиотеке libpcap и механизме фильтрации в ядре, называемом BPF \cite{mccanne1993}. Программа передаёт ядру фильтр, записанный выражением, например \texttt{tcp} или \texttt{host 10.0.1.20}, и ядро по нему отбрасывает неинтересные кадры ещё до того, как скопирует их в память приложения. Это заметно снижает нагрузку на процессор, что для одноплатного компьютера существенно. Чтобы интерфейс принимал и кадры, адресованные не ему, его нужно перевести в неразборчивый (promiscuous) режим. Поверх этого слоя работает tcpdump, который выводит захваченный трафик или сохраняет его в файл формата pcap для последующего разбора. Главное ограничение пассивного прослушивания — топология современных сетей: коммутатор доставляет кадры только на тот порт, где находится адрес назначения, поэтому анализатор захватит либо трафик, адресованный ему самому, либо трафик на зеркальном порту, либо на линии, в которую он включён напрямую. Это нужно помнить при толковании измеренной статистики.

Из захваченного трафика составляются сводные статистики. Анализатор Wireshark предлагает, в частности, разбивку трафика по протоколам, обзор общающихся станций с объёмом отправленных ими данных и распределение размеров пакетов \cite{wiresharkstat}. Для быстрой диагностики в поле полезнее всего число кадров и бит в секунду, средний размер кадра, доли отдельных протоколов и станции с наибольшим объёмом трафика: они отвечают на вопросы, общается ли в сегменте что-нибудь, какой трафик преобладает и кто его создаёт.

У самого захвата тоже есть пределы. Если программа не успевает обрабатывать кадры так быстро, как они приходят, заполняется буфер в ядре, и дальнейшие кадры ядро выбрасывает. Статистика тогда выглядит правдоподобно, но неполна. К счастью, ядро ведёт счёт выброшенных кадров, так что надёжный анализатор должен указывать это число при каждом измерении. Предел анализатора, как и у генератора, определяется скорее числом кадров в секунду, чем объёмом данных.

Для картирования сегмента снова служит Scapy: анализатор рассылает запросы ARP на все адреса выбранного диапазона и по ответам составляет список активных станций с их физическими адресами. Как отмечает Carter \cite{carter2024}, именно автоматизация таких задач на Python существенно расширяет возможности сетевого аудита, потому что позволяет подстроить собственные методики под конкретную среду.

Ещё одну группу образуют инструменты активной разведки безопасности, среди которых особое место занимает сканер Nmap. Его работа идёт в несколько этапов. Сначала он выясняет, какие адреса выбранного диапазона отвечают, затем у найденных станций исследует состояние выбранных портов и, наконец, по виду ответов пытается определить, какие службы на них слушают и какая операционная система установлена. Последняя задача, называемая снятием отпечатка стека, основана на мелких различиях в том, как разные реализации строят ответы на необычные комбинации флагов.

С активной разведкой неразрывно связан вопрос правомерности. Сканировать чужую сеть без согласия её администратора недопустимо, независимо от того, причиняет ли это вред. Проектируемое устройство решает этот вопрос своей конструкцией: оно работает в закрытом тестовом сегменте, который создаёт само и единственным другим участником которого является станция, подключённая кабелем напрямую. Сканирование поэтому никогда не выходит за пределы пространства, очерченного оператором, — это наряду с точностью измерений вторая причина логического разделения, описанного в следующем разделе.

\section{Архитектура Linux Network Namespaces (netns)}
\label{sec:netns_teorie}

Для автономного тестирования на одной физической машине используется технология Linux Network Namespaces (netns). Она обеспечивает логическую изоляцию сетевых ресурсов прямо на уровне ядра ОС. Каждое пространство имён содержит собственные независимые экземпляры сетевого стека: изолированные таблицы маршрутизации, правила межсетевого экрана и отдельные сетевые интерфейсы.

Только что созданное пространство имён полностью пусто и содержит лишь интерфейс обратной петли, к тому же выключенный. Работоспособная сетевая среда появляется только после переноса в него физического интерфейса и его поднятия. Перенесённый интерфейс одновременно исчезает из исходного пространства, так что основная система больше о нём не знает и ничего через него отправить не может.

Изоляция касается всего сетевого стека, а не только интерфейсов. У каждого пространства свои таблицы маршрутизации, своя таблица соседей, свой набор правил пакетного фильтра и свои слушающие сокеты. Два процесса в разных пространствах могут занять один и тот же порт, не мешая друг другу, что удобно для измерительного прибора с несколькими экземплярами одной службы. Команды направляются в нужное пространство вспомогательной утилитой, которая запускает процесс прямо в нём; весь механизм входит в стандартное ядро и не требует никаких надстроек.

Если пространства имён должны, наоборот, общаться, их соединяют парой виртуальных интерфейсов типа veth, которая ведёт себя как воображаемый кабель между ними. В нашем анализаторе нужно ровно противоположное, поэтому такого соединения нет. Следствие: между двумя тестовыми портами нет пути, и трафик одного не попадёт в другой даже по ошибке.

Ту же задачу можно было бы решить иначе. Можно взять два отдельных компьютера, но это противоречит требованию переносимости и повышает цену. Виртуализация добавила бы накладные расходы, которые одноплатный компьютер несёт с трудом. Остаётся маршрутизация по адресу источника или привязка приложения к конкретному интерфейсу; эти решения работают, но изоляция у них — вопрос правильной настройки, и одна ошибка или инструмент, не поддерживающий привязку к интерфейсу, её разрушит. Пространства имён, напротив, обеспечивают разделение на уровне ядра, так что неправильно направленному трафику некуда уйти, и ошибка проявляется сразу и однозначно, а не тихим искажением измеренных значений.

\section{Анализ проблематики и существующих решений}
\label{sec:analyza}

Задачу проверить состояние сетевой розетки, измерить достижимую пропускную способность или выяснить, какие устройства есть в сегменте, на практике решают тремя разными путями. У каждого свои достоинства, и именно из их сравнения вытекает задание этой работы.

Самый распространённый путь — универсальный компьютер с открытыми инструментами. Захват трафика обеспечит tcpdump или Wireshark, нагрузочные тесты — iperf3, разведку сегмента — Nmap. Это решение даёт наибольшую производительность и свободу, потому что измерительной станцией служит полноценный компьютер. Его слабость — эксплуатация. Ноутбук в шкафу или под столом использовать неудобно, для подключения к тестируемой сети ему обычно нужен переходник, и перед каждым измерением приходится вручную задавать адреса и проверять, куда на самом деле пойдёт трафик. Если же техник подключится к одному сегменту двумя интерфейсами, чтобы одновременно слушать и нагружать, он столкнётся с конфликтом таблиц маршрутизации, описанным в разделе \ref{sec:netns_teorie}: операционная система отправит трафик тем интерфейсом, который соответствует маршруту по умолчанию, а не тем, который имел в виду техник.

Второй путь — коммерческие ручные тестеры, предназначенные именно для работы в поле. Они эргономичны, сразу после включения готовы к работе и не требуют знания командной строки. Плата за удобство — закрытость. Набор функций задан производителем, добавить в прибор собственный скрипт или новую методику нельзя, а расширение обычно означает покупку старшей модели. В учебной и лабораторной среде, где важно видеть и то, как измерение проходит на самом деле, закрытая конструкция скорее мешает.

Третий путь, которому посвящена и эта работа, — одноплатный компьютер с открытой операционной системой. Пригодность платформы Raspberry Pi для этой цели подтверждена в литературе. Coşar и Kiran \cite{cosar2018} исследовали именно производительность и измеряли поведение Raspberry Pi при анализе сетевого трафика, включая загрузку процессора и памяти. Tripathi и Kumar \cite{tripathi2018} на той же платформе одновременно запускали систему обнаружения вторжений, ловушку для атакующих и анализатор пакетов, показав, что мощности платы хватает на несколько сетевых служб сразу. Hariawan и Sunaringtyas \cite{hariawan2021} построили аналогичное решение на более новом поколении платы и применили его в домашней сети.

У этих работ общая черта. На первый план в них выходит пассивная сторона задачи — захват, обнаружение и оценка трафика, который возникает в сети независимо от измерительного устройства. Активное создание нагрузки и особенно его отделение от измерительного пути в них центральной роли не играют, как и вопрос, как управлять прибором без подключённого компьютера.

Именно здесь открывается место для задания этой работы. Цель — соединить достоинства обоих крайних подходов: сохранить открытость и расширяемость компьютера с Linux и при этом достичь эргономики ручного тестера. Из сравнения существующих решений вытекают такие требования к проектируемому устройству. Прибор должен позволять одновременно пассивно слушать и активно генерировать трафик, причём обе задачи должны вестись раздельно, чтобы нагрузка не искажала захваченную статистику. Это разделение должно обеспечиваться логически, а не аккуратностью оператора, — для этого служат пространства имён. Управление должно быть возможно прямо на устройстве, без ноутбука и без физической клавиатуры, и прибор должен на время измерений работать от собственного источника энергии. IPv6 должен поддерживаться так же глубоко, как IPv4, включая автоконфигурацию адресов на другой стороне, поскольку именно она решает, сможет ли прибор обратиться к станции, на которой ничего не подготовлено. И наконец, прибор после подключения к чужой сети не должен сам ничего отправлять; активные тесты должны выполняться только по команде оператора, а нагрузочные — только в закрытом тестовом сегменте, как требует \cite{rfc6815}.

Выполнению этих требований посвящена следующая глава.
